\SourceLeaf{261}%
des Ausdrucks \(r\sin\varphi\mathord{.}r'\cos\varphi'-r\cos\varphi\mathord{.}r'\sin\varphi'\), der die Form \(=bc\sin A\) annimmt, die folgende Formel
\[
\begin{aligned}
\cos A^*-\cos A=-(q-q')p\sin A\{&\tfrac{1}{3}f^0+\tfrac{1}{6}f'p+\tfrac{1}{4}g^0(q+q')\\
&+(\tfrac{1}{10}f''-\tfrac{1}{45}f^0f^0)pp+\tfrac{3}{20}g'p(q+q')\\
&+(\tfrac{1}{5}h^0-\tfrac{7}{90}f^0f^0)(qq+qq'+q'q')+\text{etc.}\}
\end{aligned}
\]
Daraus ergibt sich weiter, bis zu Größen fünfter Ordnung,
\[
\begin{aligned}
A^*-A=-(q-q')p\{&\tfrac{1}{3}f^0+\tfrac{1}{6}f'p+\tfrac{1}{4}g^0(q+q')+\tfrac{1}{10}f''pp\\
&+\tfrac{3}{20}g'p(q+q')+\tfrac{1}{5}h^0(qq+qq'+q'q')\\
&-\tfrac{1}{90}f^0f^0(7pp+7qq+12qq'+7q'q')\}
\end{aligned}
\]
Kombinieren wir diese Formel mit der Beziehung
\[
2\sigma=ap(1-\tfrac{1}{6}f^0(pp+qq+qq'+q'q'-\text{etc.}))
\]
und mit den im vorhergehenden Artikel angegebenen Werten der Größen \(\alpha,\beta,\gamma\), so erhalten wir bis zu Größen fünfter Ordnung
\[
\begin{aligned}
\relax[11]\quad A^*=A-\sigma\{&\tfrac{1}{6}\alpha+\tfrac{1}{12}\beta+\tfrac{1}{12}\gamma+\tfrac{2}{15}f''pp+\tfrac{1}{5}g'p(q+q')\\
&+\tfrac{1}{5}h^0(3qq-2qq'+3q'q')\\
&+\tfrac{1}{90}f^0f^0(4pp-11qq+14qq'-11q'q')\}
\end{aligned}
\]
Durch ganz entsprechende Rechnungen erhalten wir
\[
\begin{aligned}
\relax[12]\quad B^*=B-\sigma\{&\tfrac{1}{12}\alpha+\tfrac{1}{6}\beta+\tfrac{1}{12}\gamma+\tfrac{1}{10}f''pp+\tfrac{1}{10}g'p(2q+q')\\
&+\tfrac{1}{5}h^0(4qq-4qq'+3q'q')\\
&-\tfrac{1}{90}f^0f^0(2pp+8qq-8qq'+11q'q')\}
\end{aligned}
\]
\[
\begin{aligned}
\relax[13]\quad C^*=C-\sigma\{&\tfrac{1}{12}\alpha+\tfrac{1}{12}\beta+\tfrac{1}{6}\gamma+\tfrac{1}{10}f''pp+\tfrac{1}{10}g'p(q+2q')\\
&+\tfrac{1}{5}h^0(3qq-4qq'+4q'q')\\
&-\tfrac{1}{90}f^0f^0(2pp+11qq-8qq'+8q'q')\}
\end{aligned}
\]
Da die Summe \(A^*+B^*+C^*\) zwei rechten Winkeln entspricht, leiten wir daraus zugleich den Überschuss der Summe \(A+B+C\) über zwei rechte Winkel ab, nämlich
\[
\begin{aligned}
\relax[14]\quad A+B+C=\pi+\sigma\{&\tfrac{1}{3}\alpha+\tfrac{1}{3}\beta+\tfrac{1}{3}\gamma+\tfrac{1}{3}f''pp+\tfrac{1}{2}g'p(q+q')\\
&+(2h^0-\tfrac{1}{3}f^0f^0)(qq-qq'+q'q')\}
\end{aligned}
\]
Diese letzte Gleichung hätte sich auch aus Formel [6] herleiten lassen.
